%%%%%%%%%%%%%%%%%%%%%%%%%%%%%%%%%%%%%%%%%%%%%%%%%%%%%%%%%%%%
% ABOUT THIS BOOK / COLOPHON
% The Composition of Advanced Mathematics
%%%%%%%%%%%%%%%%%%%%%%%%%%%%%%%%%%%%%%%%%%%%%%%%%%%%%%%%%%%%

\cleardoublepage

\chapter*{About This Book}
\label{chap:about-this-book}

\addcontentsline{toc}{chapter}{About This Book}

\markboth{About This Book}{About This Book}

%%%%%%%%%%%%%%%%%%%%%%%%%%%%%%%%%%%%%%%%%%%%%%%%%%%%%%%%%%%%
\section*{The Composition of Advanced Mathematics}
%%%%%%%%%%%%%%%%%%%%%%%%%%%%%%%%%%%%%%%%%%%%%%%%%%%%%%%%%%%%

\emph{The Composition of Advanced Mathematics} is a comprehensive,
open educational mathematics reference designed to connect the major
branches of undergraduate and graduate mathematics within a single
structured text.

Rather than treating mathematical subjects as isolated collections of
techniques, this book emphasizes the relationships among them.

The central philosophy of the text is that mathematics is built
compositionally:

\[
\boxed{
\text{foundations}
\longrightarrow
\text{structures}
\longrightarrow
\text{analysis}
\longrightarrow
\text{applications}.
}
\]

Ideas introduced in elementary mathematics reappear in increasingly
general forms throughout algebra, geometry, analysis, probability,
statistics, differential equations, optimization, computation, logic,
and mathematical modeling.

%%%%%%%%%%%%%%%%%%%%%%%%%%%%%%%%%%%%%%%%%%%%%%%%%%%%%%%%%%%%
\section*{Purpose}
%%%%%%%%%%%%%%%%%%%%%%%%%%%%%%%%%%%%%%%%%%%%%%%%%%%%%%%%%%%%

The primary goals of this book are to:

\begin{itemize}

    \item provide a broad reference across major areas of mathematics;

    \item show how mathematical subjects depend on and interact with one
          another;

    \item provide definitions, notation, theorems, formulas, examples,
          and exercises in a consistent format;

    \item support students transitioning from computational mathematics
          into proof-based mathematics;

    \item provide pathways from undergraduate mathematics into more
          advanced subjects;

    \item connect pure mathematical theory with applied mathematical
          modeling;

    \item provide a freely accessible mathematical resource that can be
          expanded and revised over time.

\end{itemize}

%%%%%%%%%%%%%%%%%%%%%%%%%%%%%%%%%%%%%%%%%%%%%%%%%%%%%%%%%%%%
\section*{Organization of the Text}
%%%%%%%%%%%%%%%%%%%%%%%%%%%%%%%%%%%%%%%%%%%%%%%%%%%%%%%%%%%%

The book is organized into fifteen major chapters.

\begin{description}

    \item[Chapter 0 --- Foundations and Mathematical Language.]
    Mathematical notation, arithmetic foundations, number systems,
    elementary set theory, relations, functions, logic, proof methods,
    infinity, and mathematical writing.

    \item[Chapter 1 --- Algebra.]
    Elementary algebra, linear algebra, abstract algebra, group theory,
    ring theory, field theory, module theory, representation theory,
    Lie theory, homological algebra, universal algebra, lattice theory,
    and category theory.

    \item[Chapter 2 --- Number Theory.]
    Divisibility, congruences, prime numbers, Diophantine equations,
    analytic number theory, algebraic number theory, computational number
    theory, elliptic curves, and arithmetic geometry.

    \item[Chapter 3 --- Geometry.]
    Euclidean, analytic, affine, projective, differential, Riemannian,
    symplectic, algebraic, and computational geometry.

    \item[Chapter 4 --- Topology.]
    Point-set topology, metric spaces, algebraic topology,
    differential topology, geometric topology, manifolds, and
    low-dimensional topology.

    \item[Chapter 5 --- Analysis.]
    Calculus, real analysis, measure theory, Lebesgue integration,
    complex analysis, functional analysis, Fourier analysis,
    harmonic analysis, operator theory, distributions,
    calculus of variations, and geometric analysis.

    \item[Chapter 6 --- Discrete Mathematics and Combinatorics.]
    Discrete structures, enumeration, recurrence relations,
    generating functions, graph theory, network flows, extremal
    combinatorics, and matroid theory.

    \item[Chapter 7 --- Probability.]
    Random variables, probability distributions, expectation,
    limit theorems, measure-theoretic probability, stochastic processes,
    stochastic calculus, stochastic differential equations,
    and random matrix theory.

    \item[Chapter 8 --- Statistics.]
    Descriptive statistics, sampling, estimation, confidence intervals,
    hypothesis testing, regression, analysis of variance,
    experimental design, Bayesian statistics, time series,
    spatial statistics, survival analysis, statistical learning,
    and computational statistics.

    \item[Chapter 9 --- Differential Equations and Dynamical Systems.]
    Ordinary differential equations, transform methods,
    dynamical systems, partial differential equations,
    difference equations, integral equations, delay equations,
    and differential-algebraic equations.

    \item[Chapter 10 --- Optimization and Operations Research.]
    Linear, nonlinear, convex, integer, combinatorial, stochastic,
    multiobjective, and network optimization together with
    dynamic programming, queueing, scheduling, inventory,
    decision theory, game theory, and optimal control.

    \item[Chapter 11 --- Numerical and Computational Mathematics.]
    Numerical analysis, numerical linear algebra, numerical optimization,
    numerical differential equations, Monte Carlo methods,
    scientific computing, symbolic computation, computer algebra,
    simulation, and high-performance mathematical computing.

    \item[Chapter 12 --- Mathematical Logic and Foundations.]
    Formal logic, axiomatic set theory, model theory, proof theory,
    computability, recursion theory, type theory,
    constructive mathematics, and philosophy of mathematics.

    \item[Chapter 13 --- Applied Mathematics and Modeling.]
    Mathematical modeling, mathematical biology, mathematical physics,
    continuum mechanics, fluid dynamics, control theory,
    network science, mathematical finance, and mathematical economics.

    \item[Chapter 14 --- Mathematical Reference.]
    Mathematical notation, symbols, number systems, constants,
    algebraic identities, geometry formulas, trigonometric identities,
    limit laws, derivative and integral tables, series, transforms,
    probability distributions, statistical tables, major theorems,
    mathematical terminology, glossary, and master index.

\end{description}

%%%%%%%%%%%%%%%%%%%%%%%%%%%%%%%%%%%%%%%%%%%%%%%%%%%%%%%%%%%%
\section*{How to Use This Book}
%%%%%%%%%%%%%%%%%%%%%%%%%%%%%%%%%%%%%%%%%%%%%%%%%%%%%%%%%%%%

This book is designed to support several styles of use.

\subsection*{Sequential Study}

A reader may begin with Chapter 0 and progress through the text in roughly
increasing mathematical sophistication.

A broad progression is

\[
\boxed{
\text{Foundations}
\rightarrow
\text{Algebra}
\rightarrow
\text{Analysis}
}
\]

followed by branches into

\[
\boxed{
\text{Geometry},
\quad
\text{Topology},
\quad
\text{Probability},
\quad
\text{Differential Equations},
\quad
\text{Optimization}.
}
\]

\subsection*{Subject Reference}

Each chapter can also be used independently as a reference for a particular
mathematical subject.

Readers may use cross-references to locate prerequisite definitions and
related ideas elsewhere in the text.

\subsection*{Formula Reference}

Chapter 14 provides compact access to frequently used formulas, identities,
symbols, distributions, and theorems.

\subsection*{Advanced Exploration}

Later sections introduce topics that may normally appear only in advanced
undergraduate or graduate courses.

These sections are intended to provide both reference material and pathways
for further study.

%%%%%%%%%%%%%%%%%%%%%%%%%%%%%%%%%%%%%%%%%%%%%%%%%%%%%%%%%%%%
\section*{Mathematical Philosophy of the Text}
%%%%%%%%%%%%%%%%%%%%%%%%%%%%%%%%%%%%%%%%%%%%%%%%%%%%%%%%%%%%

Mathematics can be viewed as a hierarchy of increasingly general structures.

For example,

\[
\boxed{
\mathbb{N}
\subset
\mathbb{Z}
\subset
\mathbb{Q}
\subset
\mathbb{R}
\subset
\mathbb{C}
}
\]

illustrates how number systems expand as new operations and problems are
considered.

Similarly,

\[
\boxed{
\text{arithmetic}
\rightarrow
\text{algebra}
\rightarrow
\text{abstract algebra}
}
\]

demonstrates how specific calculations develop into general structural
theories.

Calculus develops into analysis:

\[
\boxed{
\text{calculus}
\rightarrow
\text{real analysis}
\rightarrow
\text{measure theory}
\rightarrow
\text{functional analysis}.
}
\]

Geometry develops into increasingly structural theories:

\[
\boxed{
\text{Euclidean geometry}
\rightarrow
\text{differential geometry}
\rightarrow
\text{topology}
\rightarrow
\text{geometric analysis}.
}
\]

Applied mathematics then combines these structures:

\[
\boxed{
\text{theory}
+
\text{computation}
+
\text{data}
+
\text{modeling}.
}
\]

The title \emph{The Composition of Advanced Mathematics} reflects this
perspective: advanced mathematics is constructed through the composition
and interaction of simpler mathematical ideas.

%%%%%%%%%%%%%%%%%%%%%%%%%%%%%%%%%%%%%%%%%%%%%%%%%%%%%%%%%%%%
\section*{Notation}
%%%%%%%%%%%%%%%%%%%%%%%%%%%%%%%%%%%%%%%%%%%%%%%%%%%%%%%%%%%%

Notation is introduced when first required and consolidated in
Chapter 14.

Standard conventions are used whenever possible.

However, mathematical notation is not completely universal.

When multiple conventions are common, the text attempts to:

\begin{enumerate}

    \item state the convention being used;

    \item identify important alternative conventions;

    \item maintain that convention consistently within the relevant
          discussion.

\end{enumerate}

%%%%%%%%%%%%%%%%%%%%%%%%%%%%%%%%%%%%%%%%%%%%%%%%%%%%%%%%%%%%
\section*{Examples and Exercises}
%%%%%%%%%%%%%%%%%%%%%%%%%%%%%%%%%%%%%%%%%%%%%%%%%%%%%%%%%%%%

Worked examples are intended to demonstrate mathematical reasoning rather
than only provide final answers.

Exercises range from direct calculation to proof and theoretical
investigation.

Exercise difficulty may vary substantially because the text spans topics
from foundational mathematics through advanced undergraduate and graduate
material.

Readers are encouraged to use exercises to:

\begin{itemize}

    \item verify definitions;

    \item practice computation;

    \item construct proofs;

    \item explore counterexamples;

    \item connect multiple mathematical subjects;

    \item investigate mathematical models computationally.

\end{itemize}

%%%%%%%%%%%%%%%%%%%%%%%%%%%%%%%%%%%%%%%%%%%%%%%%%%%%%%%%%%%%
\section*{Open Educational Resource}
%%%%%%%%%%%%%%%%%%%%%%%%%%%%%%%%%%%%%%%%%%%%%%%%%%%%%%%%%%%%

This project is intended to develop as an open educational resource.

The source materials are structured so that the textbook can be revised,
expanded, and distributed in multiple formats.

The project is designed to support:

\begin{itemize}

    \item printable PDF publication;

    \item web-based publication;

    \item modular chapter distribution;

    \item accessible mathematical references;

    \item continued correction and expansion;

    \item community-based mathematical education.

\end{itemize}

%%%%%%%%%%%%%%%%%%%%%%%%%%%%%%%%%%%%%%%%%%%%%%%%%%%%%%%%%%%%
\section*{Source Formats}
%%%%%%%%%%%%%%%%%%%%%%%%%%%%%%%%%%%%%%%%%%%%%%%%%%%%%%%%%%%%

The textbook has been developed using structured mathematical source files.

The print-oriented edition uses

\[
\boxed{
\LaTeX
}
\]

for mathematical typesetting.

A web-oriented edition may also be maintained using

\[
\boxed{
\text{PreTeXt}.
}
\]

These formats allow mathematical content to be maintained separately from
the final presentation medium.

%%%%%%%%%%%%%%%%%%%%%%%%%%%%%%%%%%%%%%%%%%%%%%%%%%%%%%%%%%%%
\section*{Versioning}
%%%%%%%%%%%%%%%%%%%%%%%%%%%%%%%%%%%%%%%%%%%%%%%%%%%%%%%%%%%%

Because this is an evolving educational resource, mathematical content may
be revised as:

\begin{itemize}

    \item errors are identified;

    \item notation is standardized;

    \item examples are improved;

    \item explanations are expanded;

    \item new mathematical connections are added;

    \item references are updated;

    \item accessibility and presentation improve.

\end{itemize}

Readers should therefore consult the version or edition information
associated with the copy they are using.

%%%%%%%%%%%%%%%%%%%%%%%%%%%%%%%%%%%%%%%%%%%%%%%%%%%%%%%%%%%%
\section*{Corrections}
%%%%%%%%%%%%%%%%%%%%%%%%%%%%%%%%%%%%%%%%%%%%%%%%%%%%%%%%%%%%

No mathematical text is immune from error.

Possible issues include:

\begin{itemize}

    \item typographical errors;

    \item incorrect notation;

    \item missing hypotheses;

    \item computational mistakes;

    \item ambiguous explanations;

    \item incomplete references;

    \item formatting or cross-reference errors.

\end{itemize}

Corrections should prioritize mathematical accuracy, clarity,
and consistency across the entire text.

%%%%%%%%%%%%%%%%%%%%%%%%%%%%%%%%%%%%%%%%%%%%%%%%%%%%%%%%%%%%
\section*{Citation and Attribution}
%%%%%%%%%%%%%%%%%%%%%%%%%%%%%%%%%%%%%%%%%%%%%%%%%%%%%%%%%%%%

The development of mathematics represented in this textbook spans
centuries and draws upon the standard literature of many mathematical
disciplines.

Specific books, articles, and educational resources used in preparing
individual sections are recorded in the bibliography and source-tracking
notes maintained throughout the project.

The bibliography appears in the References section of the back matter.

%%%%%%%%%%%%%%%%%%%%%%%%%%%%%%%%%%%%%%%%%%%%%%%%%%%%%%%%%%%%
\section*{License}
%%%%%%%%%%%%%%%%%%%%%%%%%%%%%%%%%%%%%%%%%%%%%%%%%%%%%%%%%%%%

The licensing terms for distribution, modification, and reuse should be
specified for the published edition of this work.

\begin{importantbox}
Before public distribution, replace this notice with the final license
selected for the project.

For an open educational resource, a Creative Commons license such as
CC BY or CC BY-SA may be considered depending on the intended reuse
requirements.
\end{importantbox}

%%%%%%%%%%%%%%%%%%%%%%%%%%%%%%%%%%%%%%%%%%%%%%%%%%%%%%%%%%%%
\section*{Edition Information}
%%%%%%%%%%%%%%%%%%%%%%%%%%%%%%%%%%%%%%%%%%%%%%%%%%%%%%%%%%%%

\begin{center}

\textbf{\Large The Composition of Advanced Mathematics}

\vspace{0.5em}

Open Educational Mathematics Project

\vspace{1em}

\begin{tabular}{rl}
\textbf{Edition:} & Development Edition \\[0.3em]
\textbf{Format:} & \LaTeX{} / PreTeXt \\[0.3em]
\textbf{Status:} & Continuing Revision \\[0.3em]
\textbf{Language:} & English
\end{tabular}

\end{center}

%%%%%%%%%%%%%%%%%%%%%%%%%%%%%%%%%%%%%%%%%%%%%%%%%%%%%%%%%%%%
\section*{Closing Note}
%%%%%%%%%%%%%%%%%%%%%%%%%%%%%%%%%%%%%%%%%%%%%%%%%%%%%%%%%%%%

The subjects of mathematics are deeply interconnected.

Numbers lead to algebra.

Algebra produces structure.

Geometry provides space.

Topology identifies persistent form.

Analysis formalizes limits, change, and approximation.

Probability describes uncertainty.

Statistics turns observations into inference.

Differential equations describe evolution.

Optimization searches for preferred possibilities.

Computation extends what can practically be calculated.

Logic studies the foundations on which mathematical arguments are built.

Applied mathematics brings these ideas together.

Thus the structure of mathematics is not simply

\[
\boxed{
\text{a list of subjects}.
}
\]

It is better represented as

\[
\boxed{
\text{a network of mathematical ideas}.
}
\]

The purpose of this text is to make that network visible.

%%%%%%%%%%%%%%%%%%%%%%%%%%%%%%%%%%%%%%%%%%%%%%%%%%%%%%%%%%%%
% BACK MATTER SOURCE TRACKING
%%%%%%%%%%%%%%%%%%%%%%%%%%%%%%%%%%%%%%%%%%%%%%%%%%%%%%%%%%%%

%------------------------------------------------------------
% About This Book / Colophon
%------------------------------------------------------------
%
% Source basis:
%   - Internal organization and philosophy of
%     The Composition of Advanced Mathematics.
%
% Used for:
%   - project purpose;
%   - chapter overview;
%   - textbook organization;
%   - learning paths;
%   - mathematical philosophy;
%   - notation policy;
%   - examples and exercises;
%   - open educational resource statement;
%   - LaTeX and PreTeXt publication model;
%   - versioning;
%   - corrections;
%   - attribution;
%   - licensing placeholder;
%   - edition information.
%
% Notes:
%   - Replace the license placeholder before public release.
%   - Author and publication metadata can be added once finalized.
%
%%%%%%%%%%%%%%%%%%%%%%%%%%%%%%%%%%%%%%%%%%%%%%%%%%%%%%%%%%%%